\section{Method}
\label{sec:method}

\subsection{Host and provenance}

Every measurement in this paper was taken on one machine: an Apple M-series
laptop, \texttt{darwin/arm64}, 16 cores, Go 1.27.1, Node 26.5.0, on 2026-09-07.
A laptop is a deliberate choice. The claim under test is that a fleet of a
million agents does not require special hardware, and a result taken on special
hardware could not support it.

Numbers are labelled throughout. \meas{} was produced by a named script on that
host. \modeled{} was computed from measured constants and stated assumptions.
\cited{} came from a published source, with the source given. No number in this
paper is any other kind.

Each measured constant reaches the text through a macro defined in
\texttt{results-data.tex}, which is generated from the runs. If that file is
absent the paper prints \ph{} rather than a stale literal, so a number cannot
drift away from the run that produced it without the drift being visible on the
page.

\subsection{What each benchmark does}

The four scripts live in \texttt{hanzoai/cloud} under \texttt{bench/} and each
re-runs to reproduce its table.

\begin{description}[leftmargin=1.2cm,style=nextline]
\item[\texttt{fleet/fleet.mjs}] writes \rFleetAgents{} dormant agent records
  into SQLite through the production schema (100{,}000 tenants $\times$ 10
  agents), then measures the file, the write rate, and the latency of resuming
  one agent both in-process and through a cold process spawn.
\item[\texttt{goroutine/main.go}] spawns \rFleetAgents{} goroutines, each parked
  on a channel as an agent loop would be, and measures spawn latency, heap
  growth against \texttt{runtime.MemStats} after a forced GC, and the time to
  wake all of them. It also instantiates the three in-process language runtimes.
\item[\texttt{sandbox/sandbox.mjs}] measures cold start for an in-process V8
  context, a fresh OCI container, and a command dispatched into an already-warm
  container, 20 runs each.
\item[\texttt{pricing/pricing.mjs}] converts the measured volumes into a monthly
  bill at published cloud rates.
\end{description}

\subsection{Repetition, and one method note that changed the numbers}

Memory results are deterministic here and are reported as a single figure: the
per-goroutine cost came out at exactly \rGoPer{}\,bytes on all five runs.
Latency results are not deterministic, and are reported as a median with the
observed range, never as a best-of.

The one trap worth recording is that the first goroutine measurements were taken
under \texttt{go run}, which compiles the program as part of the invocation.
That inflated spawn latency by roughly a third (253\,ns against a median of
\rGoSpawn{}\,ns from a built binary) and inflated the Python context figure by
about 15\%. \emph{Benchmark a built artifact.} The rows published here come from
\texttt{go build} output, run five times, serialized so that no two benchmarks
contend for the same cores.
